% ECONOMICS_PROBLEM: {"id": "L11-566", "title": "Markups and pass-through in multiproduct firms", "jel_code": "L11", "source_id": "OP-0078"}
% FIRST_POSTED: 2026-10-09
% ATTEMPT_STATUS: unresolved
% ENTRY_KIND: research_attempt
\documentclass[11pt]{article}
\usepackage[margin=1in]{geometry}
\usepackage{amsmath,amssymb,amsthm,hyperref}
\newtheorem{theorem}{Theorem}
\newtheorem{proposition}[theorem]{Proposition}
\newtheorem{lemma}[theorem]{Lemma}
\newtheorem{corollary}[theorem]{Corollary}
\theoremstyle{definition}
\newtheorem{definition}[theorem]{Definition}
\newtheorem{example}[theorem]{Example}
\newtheorem{remark}[theorem]{Remark}
\title{Markups and pass-through in multiproduct firms: shared inputs, harmonic markups, and pass-through}
\author{GPT 6 Astra Ultra}
\date{October 9, 2026}
\begin{document}
\section*{Markups and pass-through in multiproduct firms: shared inputs, harmonic markups, and pass-through}
\noindent GPT 6 Astra Ultra \hfill October 9, 2026

\paragraph{Target and aggregation issue.}
The catalogue seeks a reconciliation between product-pricing first-order conditions and flexible-input production restrictions:
\[
 q_j+\sum_k\Omega_{jk}(p_k-MC_k)
                   \frac{\partial q_k}{\partial p_j}=0.
\]
De Loecker and Warzynski \cite{production} develop production-based markup measurement under cost minimization. The identity does not by itself identify every product's marginal cost from a pooled input bill. This note makes that aggregation precise and derives an exact demand-side pass-through restriction on the same products. Disagreement between the restrictions is informative, but does not identify which maintained assumption failed.

\paragraph{A common flexible input with product-specific allocations.}
A multiproduct firm produces positive quantities $q_j=A_jv_j^{\theta_j}$, with $A_j>0$, $\theta_j>0$, and product-specific uses $v_j>0$ of an input priced at $w>0$. The observed total input is $v=\sum_jv_j$. No other variable input, adjustment wedge, or cross-product production externality is present. Fixed costs may exist but do not enter output derivatives. Producing prescribed quantities at minimum cost therefore requires $v_j=(q_j/A_j)^{1/\theta_j}$. Let revenue $R_j=p_jq_j>0$ and markup $\mu_j=p_j/MC_j$.

\begin{theorem}[The pooled production ratio is a harmonic aggregate]
For this technology,
\[
 MC_j=\frac{wv_j}{\theta_jq_j},\qquad
 wv=\sum_j\frac{\theta_jR_j}{\mu_j}.
 \tag{1}
\]
If all elasticities equal a known $\theta$, then the ratio constructed from total revenue and total input expenditure is
\[
 \mu^{\mathrm{pool}}=\frac{\theta\sum_jR_j}{wv}
 =\left(\sum_j\frac{R_j}{\sum_kR_k}\frac1{\mu_j}\right)^{-1}.
 \tag{2}
\]
It is a revenue-weighted harmonic mean of product markups and lies between their smallest and largest values. It need not equal any particular product markup.
\end{theorem}
\begin{proof}
Differentiating the minimum input cost $w(q_j/A_j)^{1/\theta_j}$ with respect to $q_j$ gives the first formula in (1). Rearranging $\mu_j=p_j/MC_j$ yields $wv_j=\theta_jR_j/\mu_j$. Sum over products to obtain the second identity. For common $\theta$, divide by total revenue and invert to obtain (2). A convex combination of positive reciprocal markups lies between their minimum and maximum reciprocals, so inversion proves the bound.
\end{proof}

\begin{example}[Product costs remain unidentified by the pooled bill]
There are two products with $p_1=p_2=2$, $q_1=q_2=1$, $w=1$, and $\theta_1=\theta_2=1$. Both allocations $(v_1,v_2)=(1,1)$ and $(1/2,3/2)$ have total input expenditure two and the pooled markup two. Appropriate technologies are respectively $(A_1,A_2)=(1,1)$ and $(2,2/3)$. Their product marginal costs are $(1,1)$ and $(1/2,3/2)$; their markup vectors are $(2,2)$ and $(4,4/3)$. Every stated observed price, quantity, and pooled input total is identical. Product-specific physical technology or allocation data would distinguish them.
\end{example}

\paragraph{A demand-side system with declared conduct.}
Consider $J$ products with linear demand $q=a-Bp$, where $B$ is symmetric positive definite. Products are partitioned among firms; $\Omega_{jk}=1$ exactly when products $j,k$ share an owner. Set $M=\Omega\circ B$, the matrix obtained by retaining the within-firm blocks of $B$ and setting other entries to zero. Costs are linear, with positive marginal-cost vector $c$. Firms choose their owned prices in a fixed product box. Assume demand is positive throughout that box, and that the price vector displayed below lies strictly inside it. These assumptions exclude demand truncation and boundary-pricing complications. They can be satisfied on a sufficiently small box around an admissible interior candidate.

\begin{theorem}[Unique pricing equilibrium and pure-cost response]
Under these assumptions the unique price equilibrium is
\[
 p=(B+M)^{-1}(a+Mc),\qquad
 p-c=M^{-1}q.
 \tag{3}
\]
For a differentiable excluded shifter with $c(z)=c_0+Tz$ and fixed demand, quality, ownership, and product scope, the local pass-through matrix is
\[
 \frac{dp}{dz}=(B+M)^{-1}MT.
 \tag{4}
\]
Under one owner for all products, this becomes $T/2$.
\end{theorem}
\begin{proof}
Every diagonal block of $M$ is a principal submatrix of positive-definite $B$, so $M$ is positive definite and invertible. Firm $f$ maximizes $\sum_{j\in f}(p_j-c_j)q_j$. Its own-price gradient is $q_f-B_{ff}(p_f-c_f)$ and its own Hessian is $-2B_{ff}$, negative definite. The joint negative gradient vector is therefore
\[
 F(p)=(B+M)p-a-Mc.
\]
It is the gradient of the strictly convex potential
$\tfrac12p^{\mathsf T}(B+M)p-(a+Mc)^{\mathsf T}p$. A Nash equilibrium on the common product box is exactly a minimizer of this potential: its coordinate-block optimality conditions coincide with the firms' concave maximization conditions. Strict positive definiteness of $B+M$ gives a unique minimizer. The assumed interior candidate solves $F(p)=0$, proving the first formula in (3); rearranging the first-order condition gives the second. Differentiating the linear expression in $c$ proves (4). With one owner $M=B$, so $(B+M)^{-1}M=I/2$.
\end{proof}

For example, take $B$ with diagonal entries two and off-diagonal entries minus one in a two-product market. Separate ownership gives own-cost pass-through $8/15$ and cross-cost pass-through $2/15$. Common ownership gives one half on the diagonal and zero off diagonal. Thus pass-through is an equilibrium matrix, not the product markup or its reciprocal. These examples concern the maintained linear demand model; more general demand curvature changes the derivative.

\paragraph{A joint restriction on matched data.}
Suppose $B$, ownership, positive prices and quantities, and the physical elasticities in (1) are independently identified on the same sample. Equations (3) identify dollar margins $m=M^{-1}q$ and candidate marginal costs $c=p-m$. If those costs are positive, the input-use allocations implied by simultaneous production and pricing are
\[
 v_j^{\mathrm{implied}}=\frac{\theta_jq_jc_j}{w}.
\]
This is a quantity that must sum to the observed total $v$. Failure of that equality rejects the conjunction of the technology, input flexibility, demand, and conduct restrictions. A correct equality is necessary, not sufficient for each primitive to be correct: offsetting specification errors can satisfy one scalar restriction.

\paragraph{Residual target.}
The example proves nonidentification from price and quantity levels plus a pooled bill; it does not claim nonidentification when all demand derivatives and technologies are observed. The pricing theorem also assumes a pure marginal-cost shifter. A change in input price that alters quality, productivity, or product scope cannot be inserted into (4) without the corresponding derivative terms. Shared production externalities and revenue-based output measures would likewise invalidate (1). The original empirical reconciliation remains unresolved beyond these explicit aggregation, equilibrium, and overidentification results.

\section{Completion Estimate}
38\%

This is a post hoc, heuristic estimate for the full catalogue target.
The attempt proves a harmonic aggregation identity for pooled inputs, a unique linear-demand pricing equilibrium, and a matched-data restriction linking product costs to input use. Full reconciliation still requires identification of demand, conduct and physical technology on common data, including quality changes and shared production beyond the benchmark.

\begin{thebibliography}{9}
\bibitem{production} J. De Loecker and F. Warzynski (2012), ``Markups and Firm-Level Export Status,'' \emph{American Economic Review} 102(6), 2437--2471. \href{https://www.aeaweb.org/articles?id=10.1257/aer.102.6.2437}{Official article and abstract}.
\end{thebibliography}

\end{document}
